% === Begin Label Data ===
% ID: 451
% 难度星级: 
% 标签: 
% 备注: 
% 组卷引用次数: 0
% === End  Label Data ===

\begin{problem}{2015}{G}{福建卷}{15}{命题与逻辑}
一个二元码是由 $ 0 $ 和 $ 1 $ 组成的数字串 $ x_1x_2\cdots x_n $ ($ n\in\mathbb{N}^* $)，其中 $ x_k $ ($ k=1,2,\cdots,n $) 称为第 $ k $ 位码元。  
二元码是通信中常用的码，但在通信过程中有时会发生码元错误（即码元由 $ 0 $ 变为 $ 1 $，或者由 $ 1 $ 变为 $ 0 $）。  

已知某种二元码 $ x_1x_2\cdots x_7 $ 的码元满足如下校验方程组：  
$$
\begin{cases}
x_4\oplus x_5\oplus x_6\oplus x_7=0,\\
x_2\oplus x_3\oplus x_6\oplus x_7=0,\\
x_1\oplus x_3\oplus x_5\oplus x_7=0,
\end{cases}
$$  
其中运算 $ \oplus $ 定义为：$ 0\oplus 0=0 $，$ 0\oplus 1=1 $，$ 1\oplus 0=1 $，$ 1\oplus 1=0 $.  

现已知一个这种二元码在通信过程中仅在第 $ k $ 位发生码元错误后变成了 $ 1101101 $，那么利用上述校验方程组可判定 $ k $ 等于 \underline{\hspace{4em}}.
\end{problem}

\begin{answer}

\end{answer}

\begin{solutions}

\end{solutions}